\setcounter{section}{7}






  \zwischenueberschrift{Übungsaufgaben}

\inputaufgabe
{}
{





Bestimme die \anfuehrung{Kugelschnitte}{.} Welche davon sind
\definitionsverweis {Kegelschnitte}{}{?}





}
{} {}

\inputaufgabe
{}
{





Bestimme die \anfuehrung{Zylinderschnitte}{.} Welche davon sind
\definitionsverweis {Kegelschnitte}{}{?}





}
{} {}

\inputaufgabe
{}
{





Bestimme die Kegelschnitte, die sich als Schnitt mit einer Ebene ergeben, die durch den Nullpunkt verläuft.





}
{} {}

\inputaufgabe
{}
{





Zeige, dass parallele Ebenen, die beide nicht durch den Nullpunkt gehen, den gleichen Typ von Kegelschnitt definieren.





}
{} {}

\inputaufgabe
{}
{





Wir betrachten den Standardkegel

\mavergleichskette
{\vergleichskette
{ V \left(  Z^2-X^2-Y^2  \right)
}
{ \subset }{ { {\mathbb A}_{  \R  }^{  3   } }
}
{  }{
}
{    }{
}
{   }{
}
}
{}{}{}
und die Ebenen, die durch die Drehachse
\mathl{V(X,Z)}{} verlaufen. Bestimme, in Abhängigkeit vom Drehwinkel $\alpha$
\zusatzklammer {gemessen in der
\mathl{x-z}{-}Ebene gegen den Uhrzeigersinn} {} {,}
den Typ des durch die Ebene $E_\alpha$ gegebenen Kegelschnitts.





}
{} {}

\inputaufgabe
{}
{





Wir betrachten den Standardkegel

\mavergleichskette
{\vergleichskette
{ V \left(  Z^2-X^2-Y^2  \right)
}
{ \subset }{ { {\mathbb A}_{  \R  }^{  3   } }
}
{  }{
}
{    }{
}
{   }{
}
}
{}{}{}
und die Ebenen, die durch die Drehachse
\mathl{V(X-1,Z)}{} verlaufen. Bestimme, in Abhängigkeit vom Drehwinkel $\alpha$
\zusatzklammer {gemessen in der
\mathl{x-z}{-}Ebene gegen den Uhrzeigersinn} {} {,}
den Typ des durch die Ebene $E_\alpha$ gegebenen Kegelschnitts.





}
{} {}

\inputaufgabe
{}
{





Wir betrachten den Standardkegel

\mavergleichskette
{\vergleichskette
{ V \left(  Z^2-X^2-Y^2  \right)
}
{ \subset }{ { {\mathbb A}_{  \R  }^{  3   } }
}
{  }{
}
{    }{
}
{   }{
}
}
{}{}{}
und die Ebenen, die durch die Drehachse
\mathl{V(X-1,Z-1)}{} verlaufen. Bestimme, in Abhängigkeit vom Drehwinkel $\alpha$
\zusatzklammer {gemessen in der
\mathl{x-z}{-}Ebene gegen den Uhrzeigersinn} {} {,}
den Typ des durch die Ebene $E_\alpha$ gegebenen Kegelschnitts.





}
{} {}

\inputaufgabe
{}
{





Bestimme die Quadriken zu
\definitionsverweis {homogenen}{}{}
quadratischen Polynomen in zwei Variablen.





}
{} {}

\inputaufgabe
{}
{





Transformiere die Quadrik

\mathdisp {2x^2+xy-3y^2+5x-y+3} {  }

auf eine reelle Standardgestalt.





}
{} {}

\inputaufgabegibtloesung
{}
{





Betrachte die beiden Quadriken

\mathdisp {X^2+Y^2=1 \text{   und } 4X^2+3Y^2 =9} { . }

Zeige, dass die beiden Kreise über $\R$
\definitionsverweis {affin-linear äquivalent}{}{}
sind, aber nicht über $\Q$.





}
{} {Tipp: Eine Argumentationsmöglichkeit ergibt sich aus

Satz 7.2 (Maß- und Integrationstheorie (Osnabrück 2022-2023))
.}

\inputaufgabegibtloesung
{}
{





Es sei

\mavergleichskette
{\vergleichskette
{ F
}
{ = }{ (0,0)
}
{  }{
}
{    }{
}
{   }{
}
}
{}{}{}
der Nullpunkt in der reellen Ebene und

\mavergleichskette
{\vergleichskette
{ G
}
{ = }{ V(X-1)
}
{  }{
}
{    }{
}
{   }{
}
}
{}{}{.}
Es sei

\mavergleichskette
{\vergleichskette
{ e
}
{ > }{ 0
}
{  }{
}
{    }{
}
{   }{
}
}
{}{}{}
eine reelle Zahl. Bestimme eine algebraische Gleichung für die Menge der Punkte

\mavergleichskette
{\vergleichskette
{ P
}
{ = }{ (x,y)
}
{  }{
}
{    }{
}
{   }{
}
}
{}{}{}
mit der Eigenschaft, dass der Abstand $d(P,F)$ proportional mit Proportionalitätsfaktor $\sqrt{e}$ zum
\zusatzklammer {senkrechten} {} {}
Abstand $d(P,G)$ ist.

Zeige, indem Sie die Gleichung geeignet transformieren, dass bei

\mavergleichskette
{\vergleichskette
{ e
}
{ < }{ 1
}
{  }{
}
{    }{
}
{   }{
}
}
{}{}{}
eine Ellipse, bei

\mavergleichskette
{\vergleichskette
{ e
}
{ = }{ 1
}
{  }{
}
{    }{
}
{   }{
}
}
{}{}{}
eine Parabel und bei

\mavergleichskette
{\vergleichskette
{ e
}
{ > }{ 1
}
{  }{
}
{    }{
}
{   }{
}
}
{}{}{}
eine Hyperbel vorliegt.





}
{} {}

\inputaufgabe
{}
{





Parametrisiere die durch

\mavergleichskettedisp
{\vergleichskette
{F
}
{ =} { 2x^2-xy+3y^2+x-5y
}
{ } {
}
{ } {
}
{ } {
}
}
{}{}{}
definierte Quadrik mithilfe des Nullpunktes und der Geraden
\mathl{V(y-1)}{}.





}
{} {}

\inputaufgabe
{}
{





Betrachte die durch

\mavergleichskettedisp
{\vergleichskette
{ C
}
{ =} { V \left(  X^{d+1}-Y^d  \right)
}
{ } {
}
{ } {
}
{ } {
}
}
{}{}{}
definierte algebraische Kurve $C$
\zusatzklammer {
\mavergleichskettek
{\vergleichskettek
{ d
}
{ \geq }{ 1
}
{  }{
}
{    }{
}
{   }{
}
}
{}{}{}} {} {.}
Zeige, dass man mit dem Nullpunkt und der Geraden
\mathl{V(X-1)}{} eine Parametrisierung von $C$ erhält mit der im Beweis zu

Satz 7.6

beschriebenen Methode.





}
{} {}

In den folgenden Aufgaben besprechen wir ein Automorphismuskonzept, das über affin-lineare Koordinatenwechsel hinausgeht.







Es sei $K$ ein
\definitionsverweis {Körper}{}{.}
Eine
\definitionsverweis {polynomiale Abbildung}{}{}
\maabbdisp {\varphi} { { {\mathbb A}_{  K }^{  n   } } } { { {\mathbb A}_{  K }^{  n   } }
} {}
heißt
\zusatzklammer {polynomialer} {} {}
\definitionswort {Automorphismus des affinen Raumes}{,}
wenn sie eine polynomiale
\definitionsverweis {Umkehrabbildung}{}{}
besitzt.






Ein Automorphismus des affinen Raumes ist das gleiche wie ein
$K$-\definitionsverweis {Algebra-Automorphismus}{}{}
des Polynomrings
\mathl{K[X_1 , \ldots , X_n]}{} in sich. Er wird durch $n$ Polynome in $n$ Variablen gegeben.

\inputaufgabe
{}
{





Zeige, dass ein
$K$-\definitionsverweis {Algebraautomorphismus}{}{}
\maabbdisp {\varphi} {K[X] } { K[X]
} {}
durch
\mathl{X \mapsto aX+b}{} mit

\mavergleichskette
{\vergleichskette
{a
}
{ \neq }{ 0
}
{  }{
}
{    }{
}
{   }{
}
}
{}{}{}
gegeben ist
\zusatzklammer {also durch eine
\definitionsverweis {affin-lineare}{}{}
Variablentransformation} {} {.}





}
{} {}

\inputaufgabe
{}
{





Man gebe ein Beispiel für eine bijektive
\definitionsverweis {polynomiale Abbildung}{}{}
\maabbdisp {} { {\mathbb A}^{1}_{ K } } { {\mathbb A}^{1}_{ K }
} {,}
deren
\definitionsverweis {Umkehrabbildung}{}{}
nicht polynomial ist.





}
{} {}

\inputaufgabe
{}
{





Es sei

\mavergleichskette
{\vergleichskette
{ F
}
{ \in }{ K[X]
}
{  }{
}
{    }{
}
{   }{
}
}
{}{}{}
ein Polynom. Zeige, dass die Abbildung
\maabbeledisp {} { {\mathbb A}^{2}_{ K } } { {\mathbb A}^{2}_{ K }
} { (x,y) } { (x,y+F(x))
} {,}
ein
\definitionsverweis {Automorphismus}{}{}
des
\definitionsverweis {affinen Raumes}{}{}
ist. Bestimme explizit eine Umkehrabbildung.





}
{} {}

\inputaufgabe
{}
{





Bestimme die
\definitionsverweis {Umkehrabbildung}{}{}
zur Abbildung
\maabbeledisp {} { \R^2 } { \R^2
} { (x,y) } { \left(  x+y^2   , \,   -y^4-2xy^2-x^2+y^2+x+y                   \right)
} {.}





}
{} {}

\inputaufgabe
{}
{





Es sei
\maabbdisp {\varphi} { { {\mathbb A}_{  {\mathbb C}  }^{  n   } } } { { {\mathbb A}_{  {\mathbb C}  }^{  n   } }
} {}
ein
\definitionsverweis {Automorphismus}{}{}
des affinen Raumes. Zeige, dass die
\definitionsverweis {Jacobi-Determinante}{}{}
konstant gleich einem

\mavergleichskette
{\vergleichskette
{c
}
{ \neq }{ 0
}
{  }{
}
{    }{
}
{   }{
}
}
{}{}{}
ist.





}
{} {}

Das \stichwort {Jacobi-Problem} {} ist die Frage, ob eine polynomiale Abbildung
\maabbdisp {\varphi} { { {\mathbb A}_{  {\mathbb C}  }^{ n  } } } { { {\mathbb A}_{  {\mathbb C}  }^{ n  } }
} {,}
für die die Jacobi-Determinante konstant gleich $1$ ist, eine polynomiale Umkehrabbildung besitzt
\zusatzklammer {also ein Automorphismus des affinen Raumes ist} {} {.}
Dieses Problem ist schon bei

\mavergleichskette
{\vergleichskette
{ n
}
{ = }{ 2
}
{  }{
}
{    }{
}
{   }{
}
}
{}{}{}
offen. Aufgrund
des
Satzes über die Umkehrabbildung

gibt es unter der gegebenen Voraussetzung in jedem Punkt lokal eine differenzierbare Umkehrabbildung.







Zwei
\definitionsverweis {affin-algebraische Mengen}{}{}

\mavergleichskette
{\vergleichskette
{ V, \tilde{V}
}
{ \subseteq }{ { {\mathbb A}_{  K }^{  n   } }
}
{  }{
}
{    }{
}
{   }{
}
}
{}{}{}
heißen \definitionswort {affin-algebraisch äquivalent}{,} wenn es einen
\definitionsverweis {Automorphismus des affinen Raumes}{}{}
\maabb {\varphi} { { {\mathbb A}_{  K }^{  n   } }  } { { {\mathbb A}_{  K }^{  n   } }
} {}
mit

\mavergleichskettedisp
{\vergleichskette
{ \varphi^{-1}(V)
}
{ =} {\tilde{V}
}
{ } {
}
{ } {
}
{ } {
}
}
{}{}{}
gibt.







\inputaufgabe
{}
{





Es sei

\mavergleichskette
{\vergleichskette
{ F
}
{ \in }{ K[X]
}
{  }{
}
{    }{
}
{   }{
}
}
{}{}{}
ein Polynom und

\mavergleichskettedisp
{\vergleichskette
{ C
}
{ =} { V(Y-F(X))
}
{ \subset} { {\mathbb A}^{2}_{ K }
}
{ } {
}
{ } {
}
}
{}{}{}
der zugehörige
\definitionsverweis {Graph}{}{.}
Zeige, dass $C$ zur $x$-Achse

\mavergleichskette
{\vergleichskette
{ V(Y)
}
{ \subset }{ {\mathbb A}^{2}_{ K }
}
{  }{
}
{    }{
}
{   }{
}
}
{}{}{}
\definitionsverweis {affin-algebra\-isch äquivalent}{}{}
ist, aber im Allgemeinen nicht
\definitionsverweis {affin-linear äquivalent}{}{.}





}
{} {}

\inputaufgabe
{}
{





Es seien

\mavergleichskettedisp
{\vergleichskette
{V , \tilde{V}
}
{ \subset} { { {\mathbb A}_{  K  }^{  n   } }
}
{ } {
}
{ } {
}
{ } {}
}
{}{}{}
zueinander
\definitionsverweis {affin-algebraisch äquivalente}{}{}
\definitionsverweis {affin-algebraische Mengen}{}{}
mit den
\definitionsverweis {Verschwindungsidealen}{}{}
\mathkor {} {\operatorname{Id} \,(V)} {und} {\operatorname{Id} \,(\tilde{V} )} {.}
Zeige, dass dann die
\definitionsverweis {Restklassenringe}{}{}
\mathkor {} {K[X_1 , \ldots , X_n] / \operatorname{Id} \,(V)} {und} {K[X_1 , \ldots , X_n] / \operatorname{Id} \,(\tilde{V} )} {}
zueinander
\definitionsverweis {isomorph}{}{}
sind.





}
{} {}

\inputaufgabe
{}
{





Welche Quadriken im $\R^2$
\zusatzklammer {im ${\mathbb C}^2$} {} {}
sind zueinander
\definitionsverweis {affin-algebraisch äquivalent}{}{?}





}
{} {}

\inputaufgabegibtloesung
{}
{





Führe für die rationale Quadrik

\mavergleichskettedisp
{\vergleichskette
{ C
}
{ =} { V \left(  X^2+Y^2-5  \right)
}
{ \subset} { {\mathbb A}^{2}_{ \Q }
}
{ } {
}
{ } {
}
}
{}{}{}
eine rationale Parametrisierung im Sinne von

Satz 7.6

mit dem Hilfspunkt
\mathl{(1,2)}{} und einer geeigneten Geraden durch.





}
{} {}

\inputaufgabe
{}
{





Es sei $p$ eine
\definitionsverweis {Primzahl}{}{}
mit

\mavergleichskette
{\vergleichskette
{ p \mod 4
}
{ = }{ 3
}
{  }{
}
{    }{
}
{   }{
}
}
{}{}{.}
Begründe unter Bezug auf

Satz 9.10 (Zahlentheorie (Osnabrück 2025))
,
dass die rationale Quadrik

\mavergleichskettedisp
{\vergleichskette
{C
}
{ =} { V \left(  X^2+Y^2-p  \right)
}
{ \subset} { {\mathbb A}^{2}_{ \Q }
}
{ } {
}
{ } {
}
}
{}{}{}
leer ist.





}
{} {}

\inputaufgabe
{}
{





Es sei $R$ ein
\definitionsverweis {Integritätsbereich}{}{}
und

\mavergleichskette
{\vergleichskette
{r
}
{ \in }{ R
}
{  }{
}
{    }{
}
{   }{
}
}
{}{}{}
ein Element, das keine
\definitionsverweis {Quadratwurzel}{}{}
in $R$ besitze. Zeige, dass das
\definitionsverweis {Polynom}{}{}

\mavergleichskette
{\vergleichskette
{ X^2 -r
}
{ \in }{ R[X]
}
{  }{
}
{    }{
}
{   }{
}
}
{}{}{}
\definitionsverweis {irreduzibel}{}{}
ist.





}
{} {}

\inputaufgabe
{}
{





Es sei $R$ ein
\definitionsverweis {Integritätsbereich}{}{}
mit

\mavergleichskette
{\vergleichskette
{ 2
}
{ \neq }{ 0
}
{  }{
}
{    }{
}
{   }{
}
}
{}{}{}
und sei

\mavergleichskette
{\vergleichskette
{ r
}
{ \in }{ R
}
{  }{
}
{    }{
}
{   }{
}
}
{}{}{}
ein Element, das in $R$ keine
\definitionsverweis {Quadratwurzel}{}{}
besitze. Wir betrachten die
\definitionsverweis {quadratische Ringerweiterung}{}{}

\mavergleichskettedisp
{\vergleichskette
{ R
}
{ \subset} { R [X]/ \left(  X^2-r  \right)
}
{ \defeqr} { S
}
{ } {
}
{ } {
}
}
{}{}{.}
Zeige, dass die Elemente

\mavergleichskette
{\vergleichskette
{ f
}
{ \in }{ R
}
{  }{
}
{    }{
}
{   }{
}
}
{}{}{,}
die in $S$ eine Quadratwurzel besitzen, von der Form

\mavergleichskettedisp
{\vergleichskette
{ f
}
{ =} { y^2
}
{ } {
}
{ } {
}
{ } {
}
}
{}{}{}
mit

\mavergleichskette
{\vergleichskette
{ y
}
{ \in }{ R
}
{  }{
}
{    }{
}
{   }{
}
}
{}{}{}
oder von der Form

\mavergleichskettedisp
{\vergleichskette
{ f
}
{ =} { r z^2
}
{ } {
}
{ } {
}
{ } {
}
}
{}{}{}
mit

\mavergleichskette
{\vergleichskette
{ z
}
{ \in }{ R
}
{  }{
}
{    }{
}
{   }{
}
}
{}{}{}
sind.





}
{} {}

\inputaufgabe
{}
{





Es seien
\mathl{p}{} und $q$ verschiedene
\definitionsverweis {Primzahlen}{}{.}
Zeige, dass die
\zusatzklammer {über $\Q$ definierten} {} {}
Quadriken

\mavergleichskettedisp
{\vergleichskette
{C = V \left(  X^2+Y^2-p  \right) , D = V \left(  X^2+Y^2-q  \right)
}
{ \subset} { {\mathbb A}^{2}_{ \Q }
}
{ } {
}
{ } {
}
{ } {
}
}
{}{}{}
nicht zueinander
\definitionsverweis {affin-linear äquivalent}{}{}
sind.





}
{} {Tipp: Wende

Aufgabe 7.25

auf

\mavergleichskette
{\vergleichskette
{ R
}
{ = }{ K[Y]
}
{  }{
}
{    }{
}
{   }{
}
}
{}{}{}
und

\mavergleichskette
{\vergleichskette
{ S
}
{ = }{ (K[Y])[X]/ \left(  X^2-Y^2+p  \right)
}
{  }{
}
{    }{
}
{   }{
}
}
{}{}{}
an.}





  \zwischenueberschrift{Aufgaben zum Abgeben}

\inputaufgabe
{6}
{





Finde für die verschiedenen

reellen Quadriken

eine Realisierung als Kegelschnitt, also als Schnitt einer Ebene mit dem Kegel
\mathl{V(x^2+y^2-z^2)}{,} oder beweise, dass es eine solche Realisierung nicht gibt.





}
{} {}

\inputaufgabe
{9}
{





Betrachte die Abbildung
\maabbeledisp {} { {\mathbb A}^{2}_{K}  } { {\mathbb A}^{2}_{K}
} { (u,v) } { (u^2+uv,v-u^2) = (x,y)
} {.}
Bestimme zu den drei folgenden Scharen aus parallelen Geraden die Bildkurven der Geraden unter dieser Abbildung
\zusatzklammer {man gebe sowohl eine Parametrisierung als auch eine Kurvengleichung} {} {.}
Man skizziere auch typische Bildkurven für

\mavergleichskette
{\vergleichskette
{ K
}
{ = }{ \R
}
{  }{
}
{    }{
}
{   }{
}
}
{}{}{.}
\aufzaehlungdrei{Die zur $u$-Achse parallelen Geraden,
}{die zur $v$-Achse parallelen Geraden,
}{die zur Antidiagonalen parallelen Geraden.
}
Bestimme zu jeder Schar, ob sich die Bildkurven überschneiden.





}
{} {}

\inputaufgabe
{4}
{





Es seien

\mavergleichskette
{\vergleichskette
{ F,G
}
{ \in }{ K[X_1 , \ldots ,  X_n]
}
{  }{
}
{    }{
}
{   }{
}
}
{}{}{}
\definitionsverweis {Polynome}{}{}
und

\mavergleichskette
{\vergleichskette
{ K
}
{ \subseteq }{ L
}
{  }{
}
{    }{
}
{   }{
}
}
{}{}{}
eine
\definitionsverweis {Körpererweiterung}{}{.}
Diskutiere, wie sich die verschiedenen Äquivalenzbegriffe aus der siebten Vorlesung für $F$ und $G$ (und für $V(F)$ und $V(G)$) unter dem Körperwechsel verhalten.





}
{} {}

\inputaufgabe
{6}
{





Wir betrachten die beiden Restklassenringe

\mathdisp {R=\R [X,Y]/(X^2+Y^2-1) \text{   und } S=\R [X,Y]/(XY-1)} { . }

Zeige: $S$ ist ein
\definitionsverweis {Hauptidealbereich}{}{,}
$R$ hingegen nicht.





}
{\zusatzklammer {Das sind die Ringe, die zum reellen Kreis und zur reellen Hyperbel gehören.} {} {} } {Tipp: Man betrachte für $R$ das Ideal $(X-1,Y)$.}



\inputaufgabe
{4}
{





Parametrisiere die durch

\mavergleichskettedisp
{\vergleichskette
{F
}
{ =} { x^2+2xy-y^2+x-3y+4
}
{ } {
}
{ } {
}
{ } {
}
}
{}{}{}
definierte Quadrik

\mavergleichskette
{\vergleichskette
{ C
}
{ = }{ V(F)
}
{  }{
}
{    }{
}
{   }{
}
}
{}{}{}
mithilfe des Punktes

\mavergleichskette
{\vergleichskette
{ (1,2)
}
{ \in }{ C
}
{  }{
}
{    }{
}
{   }{
}
}
{}{}{}
und der $y$-Achse. Führe keine Variablentransformation durch.





}
{} {}

\inputaufgabe
{6}
{





Betrachte in ${\mathbb A}^{2}_{ \R }$ die beiden Nullstellenmengen

\mathdisp {K=V \left(  X^2+Y^2-1  \right) \text{   und } C=V \left(  X^4+Y^4-1  \right)} { . }

Zeige, dass es eine polynomiale Abbildung in zwei Variablen gibt, die die eine Nullstellenmenge surjektiv auf die andere abbildet. Zeige, dass diese Abbildung schon über $\Q$ definiert ist, dort aber nicht surjektiv ist. Zeige ferner, dass es über $\Q$ überhaupt keine surjektive polynomiale Abbildung von $C$ nach $K$ geben kann und dass es nur die konstanten polynomialen Abbildungen von $K$ nach $C$ gibt.





}
{} {}

\inputaufgabe
{6}
{





Es sei $K$ ein
\definitionsverweis {algebraisch abgeschlossener Körper}{}{}
und sei

\mavergleichskette
{\vergleichskette
{ F
}
{ \in }{ K[X,Y]
}
{  }{
}
{    }{
}
{   }{
}
}
{}{}{}
ein
\definitionsverweis {irreduzibles}{}{}
Polynom. Zeige, dass die Kurve $V(F)$ genau dann
\definitionsverweis {rational}{}{}
ist, wenn es einen injektiven
$K$-\definitionsverweis {Algebrahomomorphismus}{}{}
\maabbdisp {} { Q(K[X,Y]/(F)) } { K(T)
} {}
gibt.





}
{\zusatzklammer {Hier steht links der Quotientenkörper und rechts der rationale Funktionenkörper} {.} {}} {}
